\documentclass[11pt,letterpaper]{article}
\usepackage[margin=0.8in]{geometry}
\usepackage[T1]{fontenc}
\usepackage{lmodern}
\usepackage{amsmath,amssymb,amsthm}
\usepackage[hidelinks]{hyperref}
\usepackage{microtype}
\setlength{\parindent}{0pt}
\setlength{\parskip}{5pt}
\setlength{\abovedisplayskip}{7pt}
\setlength{\belowdisplayskip}{7pt}
\newcommand{\E}{\mathbb E}
\newcommand{\Prob}{\mathbb P}
\newcommand{\hbin}{h_2}
\newtheorem*{theorem}{Nearest-cluster entropy theorem}
\hypersetup{pdftitle={System-wide entropy from one nearest-cluster distance},pdfauthor={J. R. Landers},pdfsubject={Growing histories, neighborhood geometry, and entropy inference}}
\begin{document}
\begin{center}
{\Large\bfseries System-wide entropy from one\\[2pt] nearest-cluster distance}\par
\vspace{3pt}
{\small Growing histories, neighborhood geometry, and a shared statistical law}
\par\vspace{3pt}
{\small J. R. Landers}
\end{center}
\vspace{-5pt}

\textbf{The question.}
As observations accumulate, a history acquires new coordinates. Let
$u^{(n)}=(u_1,\ldots,u_n)$ and $v^{(n)}=(v_1,\ldots,v_n)$ be two histories
of $n$ observations, where $u_t,v_t$ are their values at coordinate $t$.
Writing $\|\cdot\|$ for the Euclidean norm, each new coordinate gives
\[
 \|u^{(n+1)}-v^{(n+1)}\|^2
 =\|u^{(n)}-v^{(n)}\|^2+(u_{n+1}-v_{n+1})^2.
\]
Histories consequently separate and fixed-radius neighborhoods can empty.
Can a measurement at just one cluster reveal the entropy of the larger system?
The answer is yes in an identifiable statistical model: the same unknown
parameter must govern both centroid separation and constituent uncertainty.
The growing record then turns one distance into many observations.

\textbf{Model and observable.}
Fix $K\geq2$ clusters, each containing $m\geq1$ labeled constituents.
Use $i=1,\ldots,K$ for clusters, $\ell=1,\ldots,m$ for constituents,
and $t=1,2,\ldots$ for coordinates. Let $B_{i\ell t}$ be mutually independent
binary variables, equal to $1$ with common unknown probability $p\in(0,1)$
and $0$ otherwise. Constituent $\ell$ of cluster $i$ has value
$X_{i\ell t}=aB_{i\ell t}$ at coordinate $t$, where $a>0$ is known.
After $n$ coordinates, define the centroid vector $C_i^{(n)}$, whose
$t$th entry is $C_{it}$, and its nearest-neighbor distance $R_n$ from cluster $1$:
\[
 C_i^{(n)}=\left(\frac1m\sum_{\ell=1}^m X_{i\ell t}\right)_{t=1}^n,
 \qquad R_n=\min_{2\leq j\leq K}\|C_1^{(n)}-C_j^{(n)}\|.
\]
Here $j$ indexes the other clusters. Only this distance is measured;
the neighbor's identity need not be retained. Write $H_n$ for the joint
Shannon entropy of all $Kmn$ constituent values, using the natural logarithm
$\ln$ (entropy in nats). For a binary probability $q\in(0,1)$, define
$\hbin(q)=-q\ln q-(1-q)\ln(1-q)$.

\begin{theorem}
With $K,m,a$ fixed as $n\to\infty$, the normalized squared distance has
limit $r$ below; $\xrightarrow{\mathrm{a.s.}}$ denotes convergence with probability one:
\begin{equation}
 \frac{R_n^2}{n}\xrightarrow{\mathrm{a.s.}}r=\frac{2a^2}{m}p(1-p),
 \qquad
 \frac{H_n}{Kmn}
 =\hbin\!\left(\frac{1-\sqrt{1-2mr/a^2}}{2}\right).
 \label{eq:main}
\end{equation}
Thus the limiting normalized distance determines the entropy per constituent
per coordinate of the whole system.
\end{theorem}

\begin{proof}
Write $\E$ for expectation and $\operatorname{Var}$ for variance.
Each centroid coordinate has mean $ap$ and variance $a^2p(1-p)/m$.
For each $j\ne1$, independence between clusters gives
\[
 \E\bigl[(C_{1t}-C_{jt})^2\bigr]
 =2\operatorname{Var}(C_{1t})=\frac{2a^2}{m}p(1-p)=r.
\]
The squared differences are independent and identically distributed across
$t$, so the strong law gives $\|C_1^{(n)}-C_j^{(n)}\|^2/n\to r$ almost
surely. The minimum of finitely many such sequences has the same limit.
Meanwhile independence of all constituents gives $H_n=Kmn\hbin(p)$.
Solving for $p$ leaves the two possibilities $p$ and $1-p$, whose entropies
are equal, yielding \eqref{eq:main}.
\end{proof}

\newpage
\textbf{An estimator from a finite record.}
Define $\widehat r_n$ as the clipped distance estimate and $\widehat h_n$
as the estimated entropy per constituent coordinate:
\begin{equation}
 \widehat r_n=\min\!\left\{\frac{a^2}{2m},\frac{R_n^2}{n}\right\},
 \qquad
 \widehat h_n=\hbin\!\left(\frac{1-\sqrt{1-2m\widehat r_n/a^2}}{2}\right).
 \label{eq:estimate}
\end{equation}
Use $\hbin(0)=\hbin(1)=0$ by continuity. Then
$\widehat h_n\to\hbin(p)$ almost surely. If the counts are known,
$\widehat H_n=Kmn\widehat h_n$ estimates the total entropy, with
\[
 \frac{\widehat H_n-H_n}{Kmn}\xrightarrow{\mathrm{a.s.}}0.
\]
This is consistency per constituent coordinate; it does not assert that
the absolute error in an increasingly large total vanishes.
For example, with $a=1$ and $m=10$, a limiting value $r=0.05$ implies
$p=1/2$ and entropy $\ln2$ per constituent coordinate. A value $r=0.018$
implies $p=0.1$ or $0.9$, both giving entropy approximately $0.3251$ nats.

For a finite-record guarantee, write $\Prob$ for probability and let
$\varepsilon>0$ be an error tolerance in normalized squared distance.
Each squared coordinate difference lies in $[0,a^2]$; Hoeffding's inequality
and a union bound over the $K-1$ candidates give
\begin{equation}
 \Prob\!\left(\left|\frac{R_n^2}{n}-r\right|>\varepsilon\right)
 \leq 2(K-1)\exp\!\left(-\frac{2n\varepsilon^2}{a^4}\right).
 \label{eq:concentration}
\end{equation}
Candidate distances share the reference cluster, but independence between
different candidates is unnecessary for this union bound.

\textbf{What the neighborhood contributes.}
Equation~\eqref{eq:main} also gives $R_n/\sqrt n\to\sqrt r>0$ almost surely.
Every fixed-radius reference neighborhood therefore eventually loses its
last other centroid. The radius needed to retain a neighbor grows as
$\sqrt{nr}$. With $a$ and $m$ fixed, its coefficient determines the entropy
density through the binary model. The nearest neighbor is useful because
even the last surviving neighbor carries the shared statistical signal.
Here the substantive information is in the calibrated growth coefficient,
rather than in the bare fact that distances grow.

\textbf{Why the shared law is essential.}
Place centroid $i$ at $c_i^{(n)}=(i,\ldots,i)$ and give each cluster two
constituents at equal and opposite coordinate offsets. In system A, every
offset magnitude is deterministically $1/4$. In system B, independently for
each cluster and coordinate, the magnitude is $1/4$ or $1/3$ with equal
probability. Both systems have identical centroids and nearest distances
$R_n=\sqrt n$. Yet their constituent entropies are respectively $0$ and
$Kn\ln2$. Centroid geometry alone cannot reveal arbitrary internal entropy.
Known scale $a$, known cluster size $m$, a common binary law, independent
coordinates, and the fixed number of clusters make the theorem identifiable.

\textbf{Thermodynamic meaning and scope.}
If the binary values encode chamber occupancy, multiplying Shannon entropy
by Boltzmann's constant $k_{\mathrm B}$ gives the coarse-grained mixing entropy.
If coordinates instead record successive independent observations,
$H_n$ is entropy of the accumulated record. It can grow at equilibrium
while the entropy of the instantaneous physical state stays constant.
The result establishes inference of system-wide uncertainty from a local
distance measurement; a derivation of thermodynamic irreversibility would
require additional dynamical information.

\vfill
{\footnotesize\textit{Project source:} \nolinkurl{expanding-neighborhoods-transcript.md},
``A theorem for your coordinate-growth model.'' The concentration estimate
above makes the finite-record implication explicit.}
\end{document}
